% from ega4-16-fr.tex
\oldpage[IV-4]{6}
système projectif de faisceaux d'anneaux sur $Y$, l'homomorphisme de transition
$\varphi_{nm}:\mathcal{O}_{Y^{(m)}}\to\mathcal{O}_{Y^{(n)}}$ pour $n\leq m$
identifiant $\mathcal{O}_{Y^{(n)}}$ au quotient de $\mathcal{O}_{Y^{(m)}}$ par la
puissance $(\mathscr{I}_f/\mathscr{I}_f^{n+1})^m$ de l'Idéal d'augmentation de
$\mathcal{O}_{Y^{(n)}}$, noyau de
$\varphi_{0n}:\mathcal{O}_{Y^{(n)}}\to\mathcal{O}_Y$. Les $Y^{(n)}$ forment par
suite un système inductif d'espaces annelés, ayant tous pour espace sous-jacent
l'espace $Y$, et on a des morphismes canoniques d'espaces annelés
$h_n:Y^{(n)}\to X$ égaux à $(\psi,\theta_n)$, où $\theta_n^\#$ est le morphisme
canonique
$\psi^*(\mathcal{O}_X)\to\psi^*(\mathcal{O}_X)/\mathscr{I}_f^{n+1}$.
Il est clair que le faisceau $\mathscr{G}r_\bullet(f)$ est un faisceau d'algèbres
graduées sur le faisceau d'anneaux
$\mathcal{O}_Y=\mathscr{G}r_0(f)$, et les $\mathscr{G}r_k(f)$ des
$\mathcal{O}_Y$-Modules.

Comme pour tout faisceau d'anneaux filtré, on a un homomorphisme canonique
surjectif de $\mathcal{O}_Y$-Algèbres graduées
\[
  \mathscr{S}_{\mathcal{O}_Y}^{\bullet}(\mathscr{G}r_1(f))
  \longrightarrow\mathscr{G}r_\bullet(f)
  \tag{16.1.2.2}
  \label{IV.16.1.2.2}
\]
coïncidant en degrés 0 et 1 avec les homomorphismes identiques.

\begin{examples}[16.1.3]
\label{IV.16.1.3}
\begin{enumerate}
\item[(i)] Supposons que $X$ soit un espace annelé en anneaux locaux, que $Y$ soit
réduit à un seul point $y$ (muni d'un anneau $\mathcal{O}_y$) et que, si
$x=\psi(y)$, $\theta^\#:\mathcal{O}_x\to\mathcal{O}_y$ soit un homomorphisme
surjectif d'anneaux ayant pour noyau l'idéal maximal $\mathfrak{m}_x$ de
$\mathcal{O}_x$. Alors les $\mathcal{O}_{Y^{(n)}}$ s'identifient aux anneaux
$\mathcal{O}_x/\mathfrak{m}_x^{n+1}$, et $\mathscr{G}r_\bullet(f)$ à l'anneau
gradué associé à l'anneau local $\mathcal{O}_x$ muni de sa filtration
$\mathfrak{m}_x$-préadique.
\item[(ii)] Supposons que $Y$ soit une partie fermée d'un sous-espace ouvert $U$
de $X$ et que $\mathcal{O}_Y$ soit induit sur $Y$ par un faisceau quotient
$\mathcal{O}_U/\mathscr{I}$, où $\mathscr{I}$ est un Idéal de $\mathcal{O}_U$
tel que $\mathscr{I}_x=\mathcal{O}_x$ pour tout $x\notin Y$; si $X$ est un
espace annelé en anneaux locaux, nous supposerons de plus que
$\mathscr{I}_x\neq\mathcal{O}_x$ pour $x\in Y$, de sorte que
$(Y,\mathcal{O}_Y)$ est encore un espace annelé en anneaux locaux.

Soit $\psi_0:Y\to U$ l'injection canonique, et notons
$\theta_0:\mathcal{O}_U\to(\psi_0)_*(\mathcal{O}_Y)$ l'homomorphisme tel que
$\theta_0^\#$ soit l'homomorphisme canonique
$\psi_0^*(\mathcal{O}_U)=\mathcal{O}_U|Y\to
(\mathcal{O}_U/\mathscr{I})|Y$, de sorte que
$j_0=(\psi_0,\theta_0):Y\to U$ est un morphisme d'espaces annelés (et d'espaces
annelés en anneaux locaux si $X$ est un espace annelé en anneaux locaux); si
$i:U\to X$ est l'injection canonique (morphisme d'espaces annelés),
$j=i\circ j_0$ est le morphisme $(\psi,\theta)$ de $Y$ dans $X$, où
$\psi:Y\to X$ est l'injection canonique, et
$\theta:\mathcal{O}_X\to\psi_*(\mathcal{O}_Y)$ est l'homomorphisme tel que
$\theta^\#=\theta_0^\#$. Comme $\theta^\#$ est surjectif, on peut appliquer les
définitions précédentes; $\mathcal{O}_{Y^{(n)}}$ est égal à
$\psi_0^*(\mathcal{O}_U/\mathscr{I}^{n+1})$, et l'on a
$(\psi_0)_*(\mathcal{O}_{Y^{(n)}})=\mathcal{O}_U/\mathscr{I}^{n+1}$, et
$\mathscr{G}r_n(j)=\mathscr{G}r_n(j_0)=
\psi_0^*(\mathscr{I}^n/\mathscr{I}^{n+1})=
j_0^*(\mathscr{I}^n/\mathscr{I}^{n+1})$.
\end{enumerate}
\end{examples}

\begin{env}[16.1.4]
\label{IV.16.1.4}
L'exemple (16.1.3, (ii)) montre qu'en général les
$\mathcal{O}_{Y^{(n)}}$ ne sont pas munis canoniquement d'une structure de
$\mathcal{O}_Y$-Module, ni a fortiori d'une structure de
$\mathcal{O}_Y$-Algèbre. La donnée d'une telle structure équivaut à celle d'un
homomorphisme de faisceaux d'anneaux
$\lambda_n:\mathcal{O}_Y\to\mathcal{O}_{Y^{(n)}}$, inverse à droite de
l'homomorphisme d'augmentation $\varphi_{0n}$; cela revient aussi à la donnée
d'un morphisme d'espaces annelés
$(1_Y,\lambda_n):Y^{(n)}\to Y$, inverse à gauche du morphisme canonique
$(1_Y,\varphi_{0n}):Y\to Y^{(n)}$.
\end{env}

\begin{proposition}[16.1.5]
\EGAReaderTarget{ega4.c04.prop.16-1-5}{IV.16.1.5}
Soit $f=(\psi,\theta):Y\to X$ une immersion de préschémas. Alors:
\begin{enumerate}
\item[(i)] $\mathscr{G}r_\bullet(f)$ est une $\mathcal{O}_Y$-Algèbre graduée
quasi-cohérente.
